\section{The \texttt{Guide\_gravity\_psd} McStas Component}
Neutron straight guide with gravity. Can be channeled and focusing.
Waviness may be specified, as well as side chamfers (on substrate).

\subsection*{Identification}
\begin{itemize}
  \item \textbf{Site:} 
  \item \textbf{Author:} \textless{}a href="mailto:farhi@ill.fr"\textgreater{}Emmanuel Farhi\textless{}/a\textgreater{}
  \item \textbf{Origin:} \textless{}a href="http://www.ill.fr"\textgreater{}ILL (France)\textless{}/a\textgreater{}.
  \item \textbf{Date:} Aug 03 2001
\end{itemize}

\subsection*{Description}
\begin{lstlisting}
Models a rectangular straight guide tube centered on the Z axis, with
gravitation handling. The entrance lies in the X-Y plane.
The guide can be channeled (nslit,d,nhslit parameters). The guide coating
specifications may be entered via different ways (global, or for
each wall m-value).

Waviness (random) may be specified either globally or for each mirror type.
Side chamfers (due to substrate processing) may be specified the same way.
In order to model a realistic straight guide assembly, a long guide of
length 'l' may be splitted into 'nelements' between which chamfers/gaps are
positioned.

The reflectivity may be specified either using an analytical mode (see
Component manual) or as a text file in free format, with 1st
column as q[Angs-1] and 2nd column as the reflectivity R in [0-1].
For details on the geometry calculation see the description in the McStas
component manual.

There is a special rotating mode in order to approximate a Fermi Chopper
behaviour, in the case the neutron trajectory is nearly linear inside the
chopper slits, i.e. the neutrons are fast w/r/ to the chopper speed.
Slits are straight, but may be super-mirror coated. In this case, the
component is NOT centered, but located at its entry window. It should then
be shifted by -l/2.

Example: Guide_gravity_psd(w1=0.1, h1=0.1, w2=0.1, h2=0.1, l=12,
R0=0.99, Qc=0.0219, alpha=6.07, m=1.0, W=0.003)

%VALIDATION
May 2005: extensive internal test, all problems solved
Validated by: K. Lieutenant
\end{lstlisting}

\subsection*{Input parameters}
Parameters in \textbf{boldface} are required; the others are optional.

\begin{longtable}{p{0.22\textwidth}p{0.12\textwidth}p{0.46\textwidth}p{0.14\textwidth}}
\toprule
\textbf{Name} & \textbf{Unit} & \textbf{Description} & \textbf{Default} \\
\midrule
\endhead
\textbf{nxpsd} & 1 & Number of bins in the built-in mirror parallel monitors &  \\
\textbf{filenameT} & str & Filename for top-mirror-monitor &  \\
\textbf{filenameB} & str & Filename for bottom-mirror-monitor &  \\
\textbf{filenameL} & str & Filename for left-mirror-monitor &  \\
\textbf{filenameR} & str & Filename for right-mirror-monitor &  \\
\textbf{w1} & m & Width at the guide entry &  \\
\textbf{h1} & m & Height at the guide entry &  \\
w2 & m & Width at the guide exit. If 0, use w1. & 0 \\
h2 & m & Height at the guide exit. If 0, use h1. & 0 \\
\textbf{l} & m & length of guide &  \\
R0 & 1 & Low-angle reflectivity & 0.995 \\
Qc & AA-1 & Critical scattering vector & 0.0218 \\
alpha & AA & Slope of reflectivity & 4.38 \\
m & 1 & m-value of material. Zero means completely absorbing. m=0.65 & 1.0 \\
W & AA-1 & Width of supermirror cut-off & 0.003 \\
nslit & 1 & Number of vertical channels in the guide (\textgreater{}= 1) (nslit-1 vertical dividing walls). & 1 \\
d & m & Thickness of subdividing walls & 0.0005 \\
mleft & 1 & m-value of material for left.   vert. mirror & -1 \\
mright & 1 & m-value of material for right.  vert. mirror & -1 \\
mtop & 1 & m-value of material for top.    horz. mirror & -1 \\
mbottom & 1 & m-value of material for bottom. horz. mirror & -1 \\
nhslit & 1 & Number of horizontal channels in the guide (\textgreater{}= 1). (nhslit-1 horizontal dividing walls). this enables to have nslit*nhslit rectangular channels & 1 \\
G & m/s2 & Gravitation norm. 0 value disables G effects. & 0 \\
aleft & 1 & alpha-value of left  vert. mirror & -1 \\
aright & 1 & alpha-value of right vert. mirror & -1 \\
atop & 1 & alpha-value of top   horz. mirror & -1 \\
abottom & 1 & alpha-value of left  horz. mirror & -1 \\
wavy & deg & Global guide waviness & 0 \\
wavy\_z & deg & Partial waviness along propagation axis & 0 \\
wavy\_tb & deg & Partial waviness in transverse direction for top/bottom mirrors & 0 \\
wavy\_lr & deg & Partial waviness in transverse direction for left/right mirrors & 0 \\
chamfers & m & Global chamfers specifications (in/out/mirror sides). & 0 \\
chamfers\_z & m & Input and output chamfers & 0 \\
chamfers\_lr & m & Chamfers on left/right mirror sides & 0 \\
chamfers\_tb & m & Chamfers on top/bottom mirror sides & 0 \\
nelements & 1 & Number of sections in the guide (length l/nelements). & 1 \\
nu & Hz & Rotation frequency (round/s) for Fermi Chopper approximation & 0 \\
phase & deg & Phase shift for the Fermi Chopper approximation & 0 \\
reflect & str & Reflectivity file name. Format q(Angs-1) R(0-1) & "NULL" \\
\bottomrule
\end{longtable}

\subsection*{Links}
\begin{itemize}
  \item \href{run:/home/willend/willend-McCode/mcstas-comps/contrib/Guide_gravity_psd.comp}{Source code} for \texttt{Guide\_gravity\_psd.comp}.
\end{itemize}
\IfFileExists{Guide_gravity_psd_static.tex}{\input{Guide_gravity_psd_static.tex}}{}